% GENERATED FILE. Edit canonical lessons, then run npm run generate:books.
% canonical-source-hash: fnv1a64:d57a2dc28144f1a8
% canonical-lessons: PT-C462-abafado, PT-C462-seca, PT-C462-aragem, PT-C462-insolacao, PT-C462-transpirar

\chapter{O calor}
\label{ch:pt-o-calor}
\hlchaptermodality{462}

\begin{chapteropening}
\textbf{By the end of this chapter:} \emph{I can talk about hot weather: muggy days, the drought, a light breeze, sunstroke and sweating.}

Complete the last lesson of chapter 462: I can talk about hot weather: muggy days, the drought, a light breeze, sunstroke and sweating.
\end{chapteropening}

\section[abafado]{abafado — muggy, stuffy, close}

\label{lesson:PT-C462-abafado}



\begin{quote}
Before the new one: say the Portuguese for the snail, then the Portuguese for to crawl.
\end{quote}

\subsection*{You'll want to know: abafado}
\textbf{abafado} — \textquotedblleft{}muggy, stuffy, close\textquotedblright{}.

\begin{quote}
Hoje o tempo está abafado.
\end{quote}

\begin{quote}
\emph{The weather is muggy today.}
\end{quote}

A room with no air is also \emph{abafado}.

\begin{cousinweb}
\emph{abafado} is built on \emph{abafar} 'to smother, to stifle'; the origin of \emph{abafar} is not settled (EtymDB lists it without a source).
\end{cousinweb}

\begin{grammarlens}[title={What you've built}]
One more word for telling, explaining and describing at B1.
\end{grammarlens}

\subsection*{Your turn}
\begin{itemize}
\raggedright
  \item \emph{Say it:} \emph{abafado}
  \item \emph{Say it:} \emph{abafado}, once more
  \item \emph{Say it:} a sentence of your own with \emph{abafado}
  \item \emph{Recall:} say the Portuguese for the snail, then the Portuguese for to crawl, then say \emph{abafado} again
\end{itemize}

\begin{tcolorbox}[breakable,colback=teal!4,colframe=teal!35!black,title={Before you move on}]
What does abafado mean? (\textquotedblleft{}Muggy, stuffy, close\textquotedblright{}.) Say it once more.
\end{tcolorbox}
\section[a seca]{a seca — the drought}

\label{lesson:PT-C462-seca}



\begin{quote}
Before the new one: say the Portuguese for to crawl, then the Portuguese for muggy.
\end{quote}

\subsection*{You'll want to know: a seca}
\textbf{a seca} — \textquotedblleft{}the drought\textquotedblright{}.

\begin{quote}
Este ano houve uma seca no sul.
\end{quote}

\begin{quote}
\emph{This year there was a drought in the south.}
\end{quote}

Informal in Portugal: \emph{Que seca!} — How boring!

\begin{cousinweb}
\emph{seca} is built on \emph{seco} 'dry', which is inherited from Latin \emph{siccus} (Lewis \& Short).
\end{cousinweb}

\begin{grammarlens}[title={What you've built}]
One more word for telling, explaining and describing at B1.
\end{grammarlens}

\subsection*{Your turn}
\begin{itemize}
\raggedright
  \item \emph{Say it:} \emph{a seca}
  \item \emph{Say it:} \emph{a seca}, once more
  \item \emph{Say it:} a sentence of your own with \emph{a seca}
  \item \emph{Recall:} say the Portuguese for to crawl, then the Portuguese for muggy, then say \emph{a seca} again
\end{itemize}

\begin{tcolorbox}[breakable,colback=teal!4,colframe=teal!35!black,title={Before you move on}]
What does a seca mean? (\textquotedblleft{}The drought\textquotedblright{}.) Say it once more.
\end{tcolorbox}
\section[a aragem]{a aragem — the light breeze}

\label{lesson:PT-C462-aragem}



\begin{quote}
Before the new one: say the Portuguese for muggy, then the Portuguese for the drought.
\end{quote}

\subsection*{You'll want to know: a aragem}
\textbf{a aragem} — \textquotedblleft{}the light breeze\textquotedblright{}.

\begin{quote}
À tarde corre uma aragem fresca.
\end{quote}

\begin{quote}
\emph{In the afternoon there is a cool breeze.}
\end{quote}

Softer than \emph{o vento}.

\begin{cousinweb}
\emph{aragem} is built on \emph{ar} 'air' with \emph{-agem} (EtymDB); \emph{ar} is inherited from Latin \emph{āēr}.
\end{cousinweb}

\begin{grammarlens}[title={What you've built}]
One more word for telling, explaining and describing at B1.
\end{grammarlens}

\subsection*{Your turn}
\begin{itemize}
\raggedright
  \item \emph{Say it:} \emph{a aragem}
  \item \emph{Say it:} \emph{a aragem}, once more
  \item \emph{Say it:} a sentence of your own with \emph{a aragem}
  \item \emph{Recall:} say the Portuguese for muggy, then the Portuguese for the drought, then say \emph{a aragem} again
\end{itemize}

\begin{tcolorbox}[breakable,colback=teal!4,colframe=teal!35!black,title={Before you move on}]
What does a aragem mean? (\textquotedblleft{}The light breeze\textquotedblright{}.) Say it once more.
\end{tcolorbox}
\section[a insolação]{a insolação — sunstroke}

\label{lesson:PT-C462-insolacao}



\begin{quote}
Before the new one: say the Portuguese for the drought, then the Portuguese for the light breeze.
\end{quote}

\subsection*{You'll want to know: a insolação}
\textbf{a insolação} — \textquotedblleft{}sunstroke\textquotedblright{}.

\begin{quote}
Ele apanhou uma insolação na praia.
\end{quote}

\begin{quote}
\emph{He got sunstroke at the beach.}
\end{quote}

\emph{Apanhar uma insolação} — to get sunstroke.

\begin{cousinweb}
\emph{insolação} is a learned borrowing from Latin \emph{īnsōlātiō} 'putting in the sun' (EtymDB), on \emph{sōl} 'sun'.
\end{cousinweb}

\begin{grammarlens}[title={What you've built}]
One more word for telling, explaining and describing at B1.
\end{grammarlens}

\subsection*{Your turn}
\begin{itemize}
\raggedright
  \item \emph{Say it:} \emph{a insolação}
  \item \emph{Say it:} \emph{a insolação}, once more
  \item \emph{Say it:} a sentence of your own with \emph{a insolação}
  \item \emph{Recall:} say the Portuguese for the drought, then the Portuguese for the light breeze, then say \emph{a insolação} again
\end{itemize}

\begin{tcolorbox}[breakable,colback=teal!4,colframe=teal!35!black,title={Before you move on}]
What does a insolação mean? (\textquotedblleft{}Sunstroke\textquotedblright{}.) Say it once more.
\end{tcolorbox}
\section[transpirar]{transpirar — to sweat, to perspire}

\label{lesson:PT-C462-transpirar}



\begin{quote}
Before the new one: say the Portuguese for the light breeze, then the Portuguese for sunstroke.
\end{quote}

\subsection*{You'll want to know: transpirar}
\textbf{transpirar} — \textquotedblleft{}to sweat, to perspire\textquotedblright{}.

\begin{quote}
Com este calor, transpiro muito.
\end{quote}

\begin{quote}
\emph{In this heat, I sweat a lot.}
\end{quote}

\emph{A transpiração} — sweat.

\begin{cousinweb}
\emph{transpirar} is a learned borrowing from Medieval Latin \emph{trānspīrāre} 'to breathe through' (EtymDB), from \emph{trāns} and \emph{spīrāre} 'to breathe'.
\end{cousinweb}

\begin{grammarlens}[title={What you've built}]
One more word for telling, explaining and describing at B1.
\end{grammarlens}

\subsection*{Your turn}
\begin{itemize}
\raggedright
  \item \emph{Say it:} \emph{transpirar}
  \item \emph{Say it:} \emph{transpirar}, once more
  \item \emph{Say it:} a sentence of your own with \emph{transpirar}
  \item \emph{Recall:} say the Portuguese for the light breeze, then the Portuguese for sunstroke, then say \emph{transpirar} again
\end{itemize}

\begin{tcolorbox}[breakable,colback=teal!4,colframe=teal!35!black,title={Before you move on}]
What does transpirar mean? (\textquotedblleft{}To sweat, to perspire\textquotedblright{}.) Say it once more.
\end{tcolorbox}
